%! AP Statistics — Unit 5 Cover Sheet
\documentclass[10pt]{article}
\usepackage{apstats-article}
\usepackage{apstats-boxes}

\begin{document}

% ── Full-bleed navy banner ────────────────────────────────────────────────────
\begin{tikzpicture}[remember picture, overlay]
  \fill[navy] (current page.north west)
    rectangle ([yshift=-1.4in]current page.north east);
\end{tikzpicture}
\vspace{-1.3in}

\begin{center}
  {\color{white}\bfseries\fontsize{28}{34}\selectfont AP Statistics}\\[8pt]
  {\color{goldacc}\bfseries\Large Shepherd \quad 2026--2027}\\[14pt]
  {\color{white}\bfseries\LARGE Unit 5: Sampling Distributions}\\[4pt]
\end{center}

\vspace{0.30in}

% ── Unit overview ─────────────────────────────────────────────────────────────
\begin{tcolorbox}[colback=sky,colframe=navy,boxrule=0.9pt,arc=2mm,
    left=4mm,right=4mm,top=3mm,bottom=3mm]
  \textbf{Unit Overview.}\quad
  Unit 5 builds the theoretical bridge between data and inference. Students begin by
  formalizing the concept of a sampling distribution --- the distribution of a statistic
  over all possible samples --- and revisiting the Normal distribution as the key model.
  The Central Limit Theorem explains why sample means cluster normally around the true
  mean even when individual values are skewed, and why larger samples produce more precise
  estimates. The unit then derives formulas for the center, spread, and shape of the
  sampling distributions of $\hat{p}$, $\hat{p}_1 - \hat{p}_2$, $\bar{x}$, and
  $\bar{x}_1 - \bar{x}_2$ --- the four statistics on which all inference in Units 6--9
  is built. Students learn to check conditions rigorously and compute probabilities for
  each statistic using the normal approximation.
\end{tcolorbox}

\vspace{0.22in}

% ── Lessons in this unit ──────────────────────────────────────────────────────
\renewcommand{\arraystretch}{1.6}
\begin{skillbox}[Lessons in This Unit]{goldbox}
\begin{tabularx}{\linewidth}{@{} c >{\bfseries}p{0.46\linewidth} X @{}}
  \textbf{\#} & \textbf{Title} & \textbf{Focus} \\
  \hline
  5.1 & Introducing Statistics: Why Is My Sample Not Like Yours?
      & Statistics vs.\ parameters; sampling variability; sampling distributions defined. \\
  \hline
  5.2 & The Normal Distribution, Revisited
      & $X \sim N(\mu,\sigma)$; forward and inverse normal probability; model appropriateness. \\
  \hline
  5.3 & The Central Limit Theorem
      & CLT conditions; $\mu_{\bar{x}} = \mu$; $\sigma_{\bar{x}} = \sigma/\sqrt{n}$; shape. \\
  \hline
  5.4 & Biased and Unbiased Point Estimates
      & Unbiased estimators ($\bar{x}$, $\hat{p}$, $s^2$); variability decreases with $n$. \\
  \hline
  5.5 & Sampling Distributions for Sample Proportions
      & $\mu_{\hat{p}} = p$; $\sigma_{\hat{p}} = \sqrt{p(1{-}p)/n}$; Large Counts condition. \\
  \hline
  5.6 & Sampling Distributions for Differences in Sample Proportions
      & $\mu = p_1{-}p_2$; $\sigma = \sqrt{p_1(1{-}p_1)/n_1 + p_2(1{-}p_2)/n_2}$; 4 counts. \\
  \hline
  5.7 & Sampling Distributions for Sample Means
      & $\mu_{\bar{x}} = \mu$; $\sigma_{\bar{x}} = \sigma/\sqrt{n}$; Normal/Large Sample. \\
  \hline
  5.8 & Sampling Distributions for Differences in Sample Means
      & $\mu = \mu_1{-}\mu_2$; $\sigma = \sqrt{\sigma_1^2/n_1 + \sigma_2^2/n_2}$; conditions. \\
  \hline
\end{tabularx}
\end{skillbox}

\vspace{0.18in}

% ── Standards covered ─────────────────────────────────────────────────────────
\begin{skillbox}[AP Statistics Big Ideas \& Learning Objectives --- Unit 5]{greenbox}
\begin{tabularx}{\linewidth}{@{} >{\bfseries}p{0.14\linewidth} X @{}}
  VAR-1.G  & Identify statistics and parameters; describe the effect of sampling variability. \\
  VAR-6.A  & Determine probabilities for a normal random variable. \\
  VAR-6.B  & Use the inverse normal to find values given probabilities. \\
  VAR-6.C  & Determine whether a normal model is appropriate for a given distribution. \\
  UNC-3.H  & Apply the Central Limit Theorem to the sampling distribution of $\bar{x}$. \\
  UNC-3.I--J & Define unbiased estimators; explain how $n$ affects variability. \\
  UNC-3.K--M & Describe center, spread, and shape for the sampling distribution of $\hat{p}$. \\
  UNC-3.N--P & Describe the sampling distribution of $\hat{p}_1 - \hat{p}_2$. \\
  UNC-3.Q--S & Describe center, spread, and shape for the sampling distribution of $\bar{x}$. \\
  UNC-3.T--V & Describe the sampling distribution of $\bar{x}_1 - \bar{x}_2$. \\
\end{tabularx}
\end{skillbox}

\end{document}
